% --- Archon physics formalization source begin ---
% archon:physics
% archon:covers PhyXMiniProblems/problem_phyx_mini_0749.lean
% archon:source-report reports/phyx_mini/problem_phyx_mini_0749.source.json

\chapter{Physics problem phyx\_mini\_0749}
\label{ch:PhyXMiniProblems_problem_phyx_mini_0749}

\paragraph{Problem source.}
\#\# Physical scenario

In figure, a sled moves in the negative \$x\$ direction at constant speed \$v\_s\$ while a ball of ice is shot from the sled with a velocity \$\textbackslash{}vec\{v\}\_0 = v\_\{0x\}\textbackslash{}hat\{i\} + v\_\{0y\}\textbackslash{}hat\{j\}\$ relative to the sled. When the ball lands, its horizontal displacement \$\textbackslash{}Delta x\_\{bg\}\$ relative to the ground (from its launch position to its landing position) is measured. Figure 4-48b gives \$\textbackslash{}Delta x\_\{bg\}\$ as a function of \$v\_s\$. Assume the ball lands at approximately its launch height.

\#\# Question

What is the value of \$v\_\{0y\}\$?

\#\# Answer choices

- A: A: 18.6 m/s

- B: B: 19.2 m/s

- C: C: 19.6 m/s

- D: D: 20.2 m/s

\#\# Figure caption (auxiliary; use the image as primary evidence)

The image consists of two parts: a diagram and a graph.

**Left Diagram:**
- Shows a sled on a horizontal surface labeled "Sled."
- A ball labeled "Ball" is positioned on the sled.
- An arrow labeled \textbackslash{}( v\_s \textbackslash{}) points to the left, indicating velocity.
- Axes are labeled \textbackslash{}( x \textbackslash{}) (horizontal) and \textbackslash{}( y \textbackslash{}) (vertical).

**Right Graph:**
- The graph plots \textbackslash{}(\textbackslash{}Delta x\_\{bg\}\textbackslash{}) (m) on the vertical axis against \textbackslash{}(v\_s\textbackslash{}) (m/s) on the horizontal axis.
- The line is straight and negatively sloped, starting at the top left (\textbackslash{}(40\textbackslash{}) on \textbackslash{}(\textbackslash{}Delta x\_\{bg\}\textbackslash{})) and ending at the bottom right (\textbackslash{}(-40\textbackslash{}) on \textbackslash{}(\textbackslash{}Delta x\_\{bg\}\textbackslash{})).
- The line crosses the horizontal axis (at \textbackslash{}(\textbackslash{}Delta x\_\{bg\} = 0\textbackslash{})) at around \textbackslash{}(v\_s = 10\textbackslash{}) m/s.
- There are grid lines for easier reading, and both axes have numeric labels.

\#\# Dataset metadata

Kinematics; Physical Model Grounding Reasoning, Multi-Formula Reasoning

\paragraph{Recorded answer/context.}
C: C: 19.6 m/s

\paragraph{Figure/image path.}
/root/proposal\_for\_physic/science-mango/phyx\_mini\_run/phyx\_data/test\_image/749.png

\paragraph{Formalization target.}
create a compiling Lean file with sorry bodies at `PhyXMiniProblems/problem\_phyx\_mini\_0749.lean`. The Lean declarations must preserve the physical quantities, dimensions or dimensional roles, figure labels, governing-law hypotheses, and final relation expressed by this problem.
Use Mathlib/Physlib names found through LeanExplore where available. If a physics API is missing, introduce faithful local abstractions rather than scalar placeholder aliases.

\begin{theorem}[Physics formalization target]
\label{thm:physics:phyx_mini_0749:target}
\lean{PhyXMiniProblems.ProblemPhyXMini0749.problem_phyx_mini_0749}
\uses{def:physics:phyx-mini-0749:phyxminiproblems-problemphyxmini0749-signedvelocityinmeterspersecond, def:physics:phyx-mini-0749:phyxminiproblems-problemphyxmini0749-sledprojectilesetup, def:physics:phyx-mini-0749:phyxminiproblems-problemphyxmini0749-matchesprimaryfigure, def:physics:phyx-mini-0749:phyxminiproblems-problemphyxmini0749-matchesproblemdescription, def:physics:phyx-mini-0749:phyxminiproblems-problemphyxmini0749-usesstandardnearearthgravity, def:physics:phyx-mini-0749:phyxminiproblems-problemphyxmini0749-hasphysicalsledprojectileparameters, def:physics:phyx-mini-0749:phyxminiproblems-problemphyxmini0749-satisfiesidealsledprojectilelaws, def:physics:phyx-mini-0749:phyxminiproblems-problemphyxmini0749-matchesdisplayedverticalspeed, def:physics:phyx-mini-0749:phyxminiproblems-problemphyxmini0749-isuniqueclosestverticalspeedchoice}
The assigned autoformalize agent should translate this physics problem into Lean declarations in the covered file, with theorem and lemma proof bodies written as `by sorry`.
\end{theorem}
\begin{proof}
This is an autoformalization task, not a proof task. Produce statements that can later be proved by the physics prover without weakening the physical model.
\end{proof}

% --- Archon declaration topology begin ---
\section{Declaration topology}
\begin{definition}[velocity Dimension]
  \label{def:physics:phyx-mini-0749:phyxminiproblems-problemphyxmini0749-velocitydimension}
  \lean{PhyXMiniProblems.ProblemPhyXMini0749.velocityDimension}
  The physical dimension of velocity, L T⁻¹
\end{definition}
\begin{proof}
  This is the declared model definition of velocity Dimension.
\end{proof}
\begin{definition}[acceleration Dimension]
  \label{def:physics:phyx-mini-0749:phyxminiproblems-problemphyxmini0749-accelerationdimension}
  \lean{PhyXMiniProblems.ProblemPhyXMini0749.accelerationDimension}
  The physical dimension of acceleration, L T⁻²
\end{definition}
\begin{proof}
  This is the declared model definition of acceleration Dimension.
\end{proof}
\begin{definition}[Signed Length Quantity]
  \label{def:physics:phyx-mini-0749:phyxminiproblems-problemphyxmini0749-signedlengthquantity}
  \lean{PhyXMiniProblems.ProblemPhyXMini0749.SignedLengthQuantity}
  A signed physical length or one-dimensional position
\end{definition}
\begin{proof}
  This is the declared model definition of Signed Length Quantity.
\end{proof}
\begin{definition}[Duration Quantity]
  \label{def:physics:phyx-mini-0749:phyxminiproblems-problemphyxmini0749-durationquantity}
  \lean{PhyXMiniProblems.ProblemPhyXMini0749.DurationQuantity}
  A nonnegative physical duration
\end{definition}
\begin{proof}
  This is the declared model definition of Duration Quantity.
\end{proof}
\begin{definition}[Signed Velocity Quantity]
  \label{def:physics:phyx-mini-0749:phyxminiproblems-problemphyxmini0749-signedvelocityquantity}
  \lean{PhyXMiniProblems.ProblemPhyXMini0749.SignedVelocityQuantity}
  \uses{def:physics:phyx-mini-0749:phyxminiproblems-problemphyxmini0749-velocitydimension}
  A signed one-dimensional physical velocity component
\end{definition}
\begin{proof}
  This is the declared model definition of Signed Velocity Quantity.
\end{proof}
\begin{definition}[Acceleration Magnitude Quantity]
  \label{def:physics:phyx-mini-0749:phyxminiproblems-problemphyxmini0749-accelerationmagnitudequantity}
  \lean{PhyXMiniProblems.ProblemPhyXMini0749.AccelerationMagnitudeQuantity}
  \uses{def:physics:phyx-mini-0749:phyxminiproblems-problemphyxmini0749-accelerationdimension}
  A nonnegative magnitude of physical acceleration
\end{definition}
\begin{proof}
  This is the declared model definition of Acceleration Magnitude Quantity.
\end{proof}
\begin{definition}[signed Length In Meters]
  \label{def:physics:phyx-mini-0749:phyxminiproblems-problemphyxmini0749-signedlengthinmeters}
  \lean{PhyXMiniProblems.ProblemPhyXMini0749.signedLengthInMeters}
  \uses{def:physics:phyx-mini-0749:phyxminiproblems-problemphyxmini0749-signedlengthquantity}
  Metre readout of a signed physical length
\end{definition}
\begin{proof}
  This is the declared model definition of signed Length In Meters.
\end{proof}
\begin{definition}[duration In Seconds]
  \label{def:physics:phyx-mini-0749:phyxminiproblems-problemphyxmini0749-durationinseconds}
  \lean{PhyXMiniProblems.ProblemPhyXMini0749.durationInSeconds}
  \uses{def:physics:phyx-mini-0749:phyxminiproblems-problemphyxmini0749-durationquantity}
  Second readout of a nonnegative physical duration
\end{definition}
\begin{proof}
  This is the declared model definition of duration In Seconds.
\end{proof}
\begin{definition}[signed Velocity In Meters Per Second]
  \label{def:physics:phyx-mini-0749:phyxminiproblems-problemphyxmini0749-signedvelocityinmeterspersecond}
  \lean{PhyXMiniProblems.ProblemPhyXMini0749.signedVelocityInMetersPerSecond}
  \uses{def:physics:phyx-mini-0749:phyxminiproblems-problemphyxmini0749-signedvelocityquantity}
  Metre-per-second readout of a signed physical velocity component
\end{definition}
\begin{proof}
  This is the declared model definition of signed Velocity In Meters Per Second.
\end{proof}
\begin{definition}[speed In Meters Per Second]
  \label{def:physics:phyx-mini-0749:phyxminiproblems-problemphyxmini0749-speedinmeterspersecond}
  \lean{PhyXMiniProblems.ProblemPhyXMini0749.speedInMetersPerSecond}
  Metre-per-second readout of a nonnegative physical speed
\end{definition}
\begin{proof}
  This is the declared model definition of speed In Meters Per Second.
\end{proof}
\begin{definition}[acceleration In Meters Per Second Squared]
  \label{def:physics:phyx-mini-0749:phyxminiproblems-problemphyxmini0749-accelerationinmeterspersecondsquared}
  \lean{PhyXMiniProblems.ProblemPhyXMini0749.accelerationInMetersPerSecondSquared}
  \uses{def:physics:phyx-mini-0749:phyxminiproblems-problemphyxmini0749-accelerationmagnitudequantity}
  Metre-per-second-squared readout of an acceleration magnitude
\end{definition}
\begin{proof}
  This is the declared model definition of acceleration In Meters Per Second Squared.
\end{proof}
\begin{definition}[speed From Meters Per Second]
  \label{def:physics:phyx-mini-0749:phyxminiproblems-problemphyxmini0749-speedfrommeterspersecond}
  \lean{PhyXMiniProblems.ProblemPhyXMini0749.speedFromMetersPerSecond}
  The physical speed having the specified nonnegative SI readout
\end{definition}
\begin{proof}
  This is the declared model definition of speed From Meters Per Second.
\end{proof}
\begin{definition}[graph Sled Speed Zero]
  \label{def:physics:phyx-mini-0749:phyxminiproblems-problemphyxmini0749-graphsledspeedzero}
  \lean{PhyXMiniProblems.ProblemPhyXMini0749.graphSledSpeedZero}
  \uses{def:physics:phyx-mini-0749:phyxminiproblems-problemphyxmini0749-speedfrommeterspersecond}
  The three sled-speed values at which the graph supplies exact points
\end{definition}
\begin{proof}
  This is the declared model definition of graph Sled Speed Zero.
\end{proof}
\begin{definition}[graph Sled Speed Ten]
  \label{def:physics:phyx-mini-0749:phyxminiproblems-problemphyxmini0749-graphsledspeedten}
  \lean{PhyXMiniProblems.ProblemPhyXMini0749.graphSledSpeedTen}
  \uses{def:physics:phyx-mini-0749:phyxminiproblems-problemphyxmini0749-speedfrommeterspersecond}
  This def records the graph Sled Speed Ten component of the problem-specific physical model
\end{definition}
\begin{proof}
  This is the declared model definition of graph Sled Speed Ten.
\end{proof}
\begin{definition}[graph Sled Speed Twenty]
  \label{def:physics:phyx-mini-0749:phyxminiproblems-problemphyxmini0749-graphsledspeedtwenty}
  \lean{PhyXMiniProblems.ProblemPhyXMini0749.graphSledSpeedTwenty}
  \uses{def:physics:phyx-mini-0749:phyxminiproblems-problemphyxmini0749-speedfrommeterspersecond}
  This def records the graph Sled Speed Twenty component of the problem-specific physical model
\end{definition}
\begin{proof}
  This is the declared model definition of graph Sled Speed Twenty.
\end{proof}
\begin{definition}[Arrow Direction]
  \label{def:physics:phyx-mini-0749:phyxminiproblems-problemphyxmini0749-arrowdirection}
  \lean{PhyXMiniProblems.ProblemPhyXMini0749.ArrowDirection}
  Directed orientations used by the coordinate axes and velocity arrow
\end{definition}
\begin{proof}
  This is the declared model definition of Arrow Direction.
\end{proof}
\begin{definition}[Figure Object]
  \label{def:physics:phyx-mini-0749:phyxminiproblems-problemphyxmini0749-figureobject}
  \lean{PhyXMiniProblems.ProblemPhyXMini0749.FigureObject}
  Objects explicitly labelled in the left-hand sled diagram
\end{definition}
\begin{proof}
  This is the declared model definition of Figure Object.
\end{proof}
\begin{definition}[Figure Quantity Label]
  \label{def:physics:phyx-mini-0749:phyxminiproblems-problemphyxmini0749-figurequantitylabel}
  \lean{PhyXMiniProblems.ProblemPhyXMini0749.FigureQuantityLabel}
  Quantity labels visible in either panel of the supplied image
\end{definition}
\begin{proof}
  This is the declared model definition of Figure Quantity Label.
\end{proof}
\begin{definition}[Graph Axis Quantity]
  \label{def:physics:phyx-mini-0749:phyxminiproblems-problemphyxmini0749-graphaxisquantity}
  \lean{PhyXMiniProblems.ProblemPhyXMini0749.GraphAxisQuantity}
  Physical quantities assigned to the two axes of the right-hand graph
\end{definition}
\begin{proof}
  This is the declared model definition of Graph Axis Quantity.
\end{proof}
\begin{definition}[Sled Projectile Figure]
  \label{def:physics:phyx-mini-0749:phyxminiproblems-problemphyxmini0749-sledprojectilefigure}
  \lean{PhyXMiniProblems.ProblemPhyXMini0749.SledProjectileFigure}
  \uses{def:physics:phyx-mini-0749:phyxminiproblems-problemphyxmini0749-arrowdirection, def:physics:phyx-mini-0749:phyxminiproblems-problemphyxmini0749-figureobject, def:physics:phyx-mini-0749:phyxminiproblems-problemphyxmini0749-figurequantitylabel, def:physics:phyx-mini-0749:phyxminiproblems-problemphyxmini0749-graphaxisquantity}
  Literal diagram and graph information from image 749.png. Schematic object coordinates express only the relative layout of the drawing; measured kinematic quantities remain in the dimensionful setup below
\end{definition}
\begin{proof}
  This is the declared model definition of Sled Projectile Figure.
\end{proof}
\begin{definition}[Sled Projectile Setup]
  \label{def:physics:phyx-mini-0749:phyxminiproblems-problemphyxmini0749-sledprojectilesetup}
  \lean{PhyXMiniProblems.ProblemPhyXMini0749.SledProjectileSetup}
  \uses{def:physics:phyx-mini-0749:phyxminiproblems-problemphyxmini0749-signedlengthquantity, def:physics:phyx-mini-0749:phyxminiproblems-problemphyxmini0749-durationquantity, def:physics:phyx-mini-0749:phyxminiproblems-problemphyxmini0749-signedvelocityquantity, def:physics:phyx-mini-0749:phyxminiproblems-problemphyxmini0749-accelerationmagnitudequantity, def:physics:phyx-mini-0749:phyxminiproblems-problemphyxmini0749-sledprojectilefigure}
  Independent physical quantities in the family of launch experiments shown by the graph. The parameter of each function is the nonnegative speed magnitude v\_s; the corresponding sled velocity component is negative. In particular, the requested vertical launch component is an independent field and is not defined from the recorded answer
\end{definition}
\begin{proof}
  This is the declared model definition of Sled Projectile Setup.
\end{proof}
\begin{definition}[Matches Primary Figure]
  \label{def:physics:phyx-mini-0749:phyxminiproblems-problemphyxmini0749-matchesprimaryfigure}
  \lean{PhyXMiniProblems.ProblemPhyXMini0749.MatchesPrimaryFigure}
  \uses{def:physics:phyx-mini-0749:phyxminiproblems-problemphyxmini0749-signedlengthinmeters, def:physics:phyx-mini-0749:phyxminiproblems-problemphyxmini0749-graphsledspeedzero, def:physics:phyx-mini-0749:phyxminiproblems-problemphyxmini0749-graphsledspeedten, def:physics:phyx-mini-0749:phyxminiproblems-problemphyxmini0749-graphsledspeedtwenty, def:physics:phyx-mini-0749:phyxminiproblems-problemphyxmini0749-sledprojectilesetup}
  Primary-image evidence. Besides the coordinate directions, object labels, and units, it records the three collinear graph points (0, 40), (10, 0), and (20, -40). These are measured inputs, not the requested value of v0y
\end{definition}
\begin{proof}
  This is the declared model definition of Matches Primary Figure.
\end{proof}
\begin{definition}[Matches Problem Description]
  \label{def:physics:phyx-mini-0749:phyxminiproblems-problemphyxmini0749-matchesproblemdescription}
  \lean{PhyXMiniProblems.ProblemPhyXMini0749.MatchesProblemDescription}
  \uses{def:physics:phyx-mini-0749:phyxminiproblems-problemphyxmini0749-sledprojectilesetup}
  The prose specifies a level launch-and-landing experiment and identifies the two stored launch-velocity components as components relative to the sled. Equal endpoint heights is scenario data; it does not prescribe either the flight duration or the requested vertical velocity
\end{definition}
\begin{proof}
  This is the declared model definition of Matches Problem Description.
\end{proof}
\begin{definition}[Uses Standard Near Earth Gravity]
  \label{def:physics:phyx-mini-0749:phyxminiproblems-problemphyxmini0749-usesstandardnearearthgravity}
  \lean{PhyXMiniProblems.ProblemPhyXMini0749.UsesStandardNearEarthGravity}
  \uses{def:physics:phyx-mini-0749:phyxminiproblems-problemphyxmini0749-accelerationinmeterspersecondsquared, def:physics:phyx-mini-0749:phyxminiproblems-problemphyxmini0749-sledprojectilesetup}
  Standard near-Earth gravitational acceleration used by the ideal model
\end{definition}
\begin{proof}
  This is the declared model definition of Uses Standard Near Earth Gravity.
\end{proof}
\begin{definition}[Has Physical Sled Projectile Parameters]
  \label{def:physics:phyx-mini-0749:phyxminiproblems-problemphyxmini0749-hasphysicalsledprojectileparameters}
  \lean{PhyXMiniProblems.ProblemPhyXMini0749.HasPhysicalSledProjectileParameters}
  \uses{def:physics:phyx-mini-0749:phyxminiproblems-problemphyxmini0749-durationinseconds, def:physics:phyx-mini-0749:phyxminiproblems-problemphyxmini0749-accelerationinmeterspersecondsquared, def:physics:phyx-mini-0749:phyxminiproblems-problemphyxmini0749-sledprojectilesetup}
  Positivity conditions selecting the nondegenerate physical flight
\end{definition}
\begin{proof}
  This is the declared model definition of Has Physical Sled Projectile Parameters.
\end{proof}
\begin{definition}[Satisfies Ideal Sled Projectile Laws]
  \label{def:physics:phyx-mini-0749:phyxminiproblems-problemphyxmini0749-satisfiesidealsledprojectilelaws}
  \lean{PhyXMiniProblems.ProblemPhyXMini0749.SatisfiesIdealSledProjectileLaws}
  \uses{def:physics:phyx-mini-0749:phyxminiproblems-problemphyxmini0749-signedlengthinmeters, def:physics:phyx-mini-0749:phyxminiproblems-problemphyxmini0749-durationinseconds, def:physics:phyx-mini-0749:phyxminiproblems-problemphyxmini0749-signedvelocityinmeterspersecond, def:physics:phyx-mini-0749:phyxminiproblems-problemphyxmini0749-speedinmeterspersecond, def:physics:phyx-mini-0749:phyxminiproblems-problemphyxmini0749-accelerationinmeterspersecondsquared, def:physics:phyx-mini-0749:phyxminiproblems-problemphyxmini0749-sledprojectilesetup}
  Ideal no-drag kinematics for every plotted value of v\_s: the sled's ground-frame x velocity is -v\_s; Galilean velocity addition gives the ball's ground-frame x velocity; horizontal displacement is that constant velocity times the common flight duration; and vertical displacement under uniform downward gravity is v0y * T - g * T\textasciicircum{}2 / 2. These are general physical laws and contain neither the numerical value of v0y nor an answer-choice assertion
\end{definition}
\begin{proof}
  This is the declared model definition of Satisfies Ideal Sled Projectile Laws.
\end{proof}
\begin{lemma}[graph Determines Horizontal Launch Velocity And Flight Duration]
  \label{lem:physics:phyx-mini-0749:phyxminiproblems-problemphyxmini0749-graphdetermineshorizontallaunchvelocityandflightduration}
  \lean{PhyXMiniProblems.ProblemPhyXMini0749.graphDeterminesHorizontalLaunchVelocityAndFlightDuration}
  \uses{def:physics:phyx-mini-0749:phyxminiproblems-problemphyxmini0749-durationinseconds, def:physics:phyx-mini-0749:phyxminiproblems-problemphyxmini0749-signedvelocityinmeterspersecond, def:physics:phyx-mini-0749:phyxminiproblems-problemphyxmini0749-sledprojectilesetup, def:physics:phyx-mini-0749:phyxminiproblems-problemphyxmini0749-matchesprimaryfigure, def:physics:phyx-mini-0749:phyxminiproblems-problemphyxmini0749-hasphysicalsledprojectileparameters, def:physics:phyx-mini-0749:phyxminiproblems-problemphyxmini0749-satisfiesidealsledprojectilelaws}
  The graph's horizontal intercept and vertical intercept, together with Galilean horizontal motion, determine v0x = 10 m/s and a 4 s flight. Neither result is included in a premise structure
\end{lemma}
\begin{proof}
  The conclusion follows from the governing relations and figure readouts named in the statement of graph Determines Horizontal Launch Velocity And Flight Duration.
\end{proof}
\begin{lemma}[vertical Launch Velocity In Meters Per Second eq 19 6]
  \label{lem:physics:phyx-mini-0749:phyxminiproblems-problemphyxmini0749-verticallaunchvelocityinmeterspersecond-eq-19-6}
  \lean{PhyXMiniProblems.ProblemPhyXMini0749.verticalLaunchVelocityInMetersPerSecond_eq_19_6}
  \uses{def:physics:phyx-mini-0749:phyxminiproblems-problemphyxmini0749-signedvelocityinmeterspersecond, def:physics:phyx-mini-0749:phyxminiproblems-problemphyxmini0749-sledprojectilesetup, def:physics:phyx-mini-0749:phyxminiproblems-problemphyxmini0749-matchesprimaryfigure, def:physics:phyx-mini-0749:phyxminiproblems-problemphyxmini0749-matchesproblemdescription, def:physics:phyx-mini-0749:phyxminiproblems-problemphyxmini0749-usesstandardnearearthgravity, def:physics:phyx-mini-0749:phyxminiproblems-problemphyxmini0749-hasphysicalsledprojectileparameters, def:physics:phyx-mini-0749:phyxminiproblems-problemphyxmini0749-satisfiesidealsledprojectilelaws}
  Equal launch and landing heights give 0 = v0y T - g T²/2; using the derived T = 4 s and g = 9.80 m/s² yields v0y = 19.6 m/s
\end{lemma}
\begin{proof}
  The conclusion follows from the governing relations and figure readouts named in the statement of vertical Launch Velocity In Meters Per Second eq 19 6.
\end{proof}
\begin{definition}[Answer Choice]
  \label{def:physics:phyx-mini-0749:phyxminiproblems-problemphyxmini0749-answerchoice}
  \lean{PhyXMiniProblems.ProblemPhyXMini0749.AnswerChoice}
  Labels of the four speed values printed with the problem
\end{definition}
\begin{proof}
  This is the declared model definition of Answer Choice.
\end{proof}
\begin{definition}[vertical Speed Meters Per Second]
  \label{def:physics:phyx-mini-0749:phyxminiproblems-problemphyxmini0749-answerchoice-verticalspeedmeterspersecond}
  \lean{PhyXMiniProblems.ProblemPhyXMini0749.AnswerChoice.verticalSpeedMetersPerSecond}
  \uses{def:physics:phyx-mini-0749:phyxminiproblems-problemphyxmini0749-answerchoice}
  Vertical-speed readout in metres per second printed beside each choice
\end{definition}
\begin{proof}
  This is the declared model definition of vertical Speed Meters Per Second.
\end{proof}
\begin{definition}[recorded Answer Choice]
  \label{def:physics:phyx-mini-0749:phyxminiproblems-problemphyxmini0749-recordedanswerchoice}
  \lean{PhyXMiniProblems.ProblemPhyXMini0749.recordedAnswerChoice}
  \uses{def:physics:phyx-mini-0749:phyxminiproblems-problemphyxmini0749-answerchoice}
  The answer label recorded by the supplied dataset metadata
\end{definition}
\begin{proof}
  This is the declared model definition of recorded Answer Choice.
\end{proof}
\begin{definition}[Matches Displayed Vertical Speed]
  \label{def:physics:phyx-mini-0749:phyxminiproblems-problemphyxmini0749-matchesdisplayedverticalspeed}
  \lean{PhyXMiniProblems.ProblemPhyXMini0749.MatchesDisplayedVerticalSpeed}
  \uses{def:physics:phyx-mini-0749:phyxminiproblems-problemphyxmini0749-signedvelocityinmeterspersecond, def:physics:phyx-mini-0749:phyxminiproblems-problemphyxmini0749-sledprojectilesetup, def:physics:phyx-mini-0749:phyxminiproblems-problemphyxmini0749-answerchoice}
  The modeled vertical launch component equals the value printed by choice
\end{definition}
\begin{proof}
  This is the declared model definition of Matches Displayed Vertical Speed.
\end{proof}
\begin{definition}[Is Unique Closest Vertical Speed Choice]
  \label{def:physics:phyx-mini-0749:phyxminiproblems-problemphyxmini0749-isuniqueclosestverticalspeedchoice}
  \lean{PhyXMiniProblems.ProblemPhyXMini0749.IsUniqueClosestVerticalSpeedChoice}
  \uses{def:physics:phyx-mini-0749:phyxminiproblems-problemphyxmini0749-signedvelocityinmeterspersecond, def:physics:phyx-mini-0749:phyxminiproblems-problemphyxmini0749-sledprojectilesetup, def:physics:phyx-mini-0749:phyxminiproblems-problemphyxmini0749-answerchoice}
  A displayed value is strictly closer to the modeled readout than every alternative
\end{definition}
\begin{proof}
  This is the declared model definition of Is Unique Closest Vertical Speed Choice.
\end{proof}
% --- Archon declaration topology end ---
% --- Archon physics formalization source end ---
